%\begin{figure}[ht]
%\centering
%    \begin{subfigure}[b]{0.9\linewidth}
%        \includegraphics[width=\linewidth]{figures/data_diagram.png}
%        \caption{Data processing and augmentation.}
%    \end{subfigure}
%
%    \begin{subfigure}[b]{0.55\linewidth}
%        \includegraphics[width=\linewidth]{figures/training_diagram.png}
%        \caption{Model training.}
%    \end{subfigure}
\begin{figure}[h]
\centering
\includegraphics[width=0.99\linewidth]{figures/acid_training.pdf}
\caption{Data processing and model training. (a) Each document-query pair from the training corpus will be converted into inputs and outputs for finetuning the pretrained transformer decoder, which serves as the generative retrieval model. (b) GPT-3.5 is used to generate a sequence of keyphrases, which is used as the document ID. (c) Given a user query or a synthetic query, the generative retrieval model learns to generate the ID of the relevant document. We use a doc2query model to generate synthetic queries as additional inputs.
Randomly sampled spans of 64 tokens can also be used as inputs to ensure that the model associates the contents of each document with its ID.
} 
\label{fig:architecture}
\end{figure}